\documentclass[a4paper, 12pt]{article}

\usepackage{amsmath, amssymb, amsthm}
\usepackage{enumerate}
\usepackage[mathscr]{eucal}
\usepackage{fancyhdr}
\usepackage[top=1.5in, bottom=1.5in, left=1.0in, right=1.0in]{geometry}
\usepackage{lastpage}
\usepackage{latexsym}
\usepackage{listings}
\usepackage{tikz}

\newcommand{\Author}[1]{\def \AuthorVar {#1}}
\newcommand{\Course}[1]{\def \CourseVar {#1}}
\newcommand{\Title}[1]{\def \TitleVar {#1}}
\newcommand{\Semester}[1]{\def \SemesterVar {#1}}

\pagestyle{fancy}
\lhead{\textbf{\CourseVar}}
\chead{\textbf{\TitleVar}}
\rhead{\textbf{\AuthorVar}}
\lfoot{\textbf{\SemesterVar}}
\cfoot{}
\rfoot{\textbf{\thepage\ of \pageref{LastPage}}}
\renewcommand{\headrulewidth}{0.4pt}
\renewcommand{\footrulewidth}{0.4pt}

\newcommand{\tbf}{\textbf}
\newcommand{\tit}{\textit}
\newcommand{\tsl}{\textsl}
\newcommand{\tsc}{\textsc}
\newcommand{\ttt}{\texttt}
\newcommand{\trm}{\textrm}
\newcommand{\tsf}{\textsf}
\newcommand{\tup}{\textup}
\newcommand{\mbf}{\mathbf}
\providecommand{\mit}{\mathit}
\newcommand{\msf}{\mathsf}
\newcommand{\mrm}{\mathm}
\newcommand{\mtt}{\mathtt}
\newcommand{\Bold}[1]{\boldsymbol{#1}}
\newcommand{\Cal}[1]{\mathcal{#1}}
\newcommand{\DD}[1]{\mathbb{#1}}
\newcommand{\Euler}[1]{\mathscr{#1}}
\newcommand{\Frak}[1]{\mathfrak{#1}}

\theoremstyle{remark}
\newtheorem{problem}{\normalfont \bfseries Problem}

\definecolor{lightgray}{RGB}{240,240,240}
\definecolor{darkgreen}{RGB}{51,102,0}

\lstset{
  backgroundcolor=\color{lightgray},
  basicstyle=\footnotesize\ttfamily,
  showstringspaces=false,
  commentstyle=\color{darkgreen},
  keywordstyle=\color{blue},
  stringstyle=\color{orange},
}

\Course{CS110}
\Title{Project 3}
\Author{Swami Iyer}
\Semester{Spring 2015}

\begin{document}

\noindent \tbf{Section Readings:} 1.5, 1.6, 2.1, 2.2

\bigskip

\begin{problem}
(\tit{Dragon Curves}) Write a program \lstinline$Dragon$ that takes an \lstinline$int$ value $N$ as command-line argument and prints the instructions for drawing the dragon curve of order $N$. The instructions are strings of F, L, and R characters, where F means ``draw line while moving 1 unit forward,'' L means ``turn left,''
and R means ``turn right.'' A dragon curve of order 0 is just the character F, and a curve of order $N$ is a curve of order $N-1$ followed by an L followed by a curve of order $N-1$ traversed in reverse order with R replaced by L and L replaced by R. 

\begin{lstlisting}[language=bash]
$ java Dragon 0
F
$ java Dragon 1
FLF
$ java Dragon 2
FLFLFRF
$ java Dragon 3
FLFLFRFLFLFRFRF
\end{lstlisting}
Dragon curves of orders 0, 1, 2, and 3 and the corresponding instructions are shown below.

\begin{center}
\includegraphics[scale=0.7]{dragon.pdf}
\end{center}
\end{problem}

\begin{problem}
(\tit{Visualizing Dragon Curves}) Write a program \lstinline$DragonPlot$ that reads from standard input the instructions produced by \lstinline$Dragon$ (from Problem 1) for drawing a dragon curve, generates a drawing of the curve (shown below), and saves it in a file called \lstinline$dragon.jpg$. 

\begin{lstlisting}[language=bash]
$ java Dragon 13 | java DragonPlot
\end{lstlisting}

\begin{center}
\includegraphics[scale=0.35]{dragon_reference.jpg}
\end{center}
\end{problem}

\begin{problem}
(\tit{Numbers Library}) Implement \lstinline$Numbers$, a library of static methods with the following API:

\begin{lstlisting}[language=Java]
public class Numbers

    // Return true if n is prime, and false otherwise.
    public static boolean isPrime(int n)
    
    // Return the greatest common divisor of a and b.
    public static int gcd(int a, int b)
    
    // Return true if a and b are coprime, ie, their gcd is 1, and false 
    // otherwise.
    public static boolean coprime(int a, int b)

    // Return the sum of the proper divisors of n. Eg, if n is 6, return 
    // 1 + 2 + 3 = 6, since 1, 2, and 3 are the proper divisors of 6.
    public static int sumOfProperDivisors(int n) 

    // Return true if n is perfect, ie, its proper divisors add up to n, 
    // and false otherwise.
    public static boolean isPerfect(int n) 
    
    // Return true if a and b are amicable, ie, the proper divisors of a 
    // add up to b and the proper divisors of b add up to a.
    public static boolean amicable(int a, int b) 
\end{lstlisting}

\begin{lstlisting}[language=bash]
$ java Numbers
false
true
8
true
true
false
true
\end{lstlisting}
\end{problem}

\begin{problem}
(\tit{Counting Primes}) Write a client program \lstinline$PrimeCounter$ that takes an \lstinline$int$ value $n$ as command-line argument and prints the number of primes less than or equal to $n$. Your implementation must use the \lstinline$isPrime()$ method from \lstinline$Numbers$ to check if a number is prime or not.

\begin{lstlisting}[language=bash]
$ java PrimeCounter 100
25
$ java PrimeCounter 10000000
664579
\end{lstlisting}
\end{problem}

\begin{problem}
(\tit{Counting Totatives}) Write a client program \lstinline$TotativeCounter$ that takes an \lstinline$int$ value $n$ as command-line argument and prints the number totatives of $n$, ie, the number of positive integers less than or equal to $n$ that are coprime to $n$. Your implementation must use the \lstinline$coprime()$ method from \lstinline$Numbers$ to check if two numbers are coprime.

\begin{lstlisting}[language=bash]
$ java TotativeCounter 100
40
\end{lstlisting}
\end{problem}

\begin{problem}
(\tit{Coprime Pattern}) Write a client program \lstinline$CoprimePattern$ that takes an \lstinline$int$ value $n$ as a command-line argument and prints an $n$-by-$n$ matrix such that the
element in row $i$ and column $j$ $(1 \leq i, j \leq n)$ is a \lstinline$"*"$ (a star) if $i$ and $j$ are coprime and a \lstinline$" "$ (a space) otherwise. The row numbers should be
printed at the end of each row. Your implementation must use the \lstinline$coprime()$ method from \lstinline$Numbers$ to check if two numbers are coprime.

\begin{lstlisting}[language=bash]
$ java CoprimePattern 10
* * * * * * * * * * 1
*   *   *   *   *   2
* *   * *   * *   * 3
*   *   *   *   *   4
* * * *   * * * *   5
*       *   *       6
* * * * * *   * * * 7
*   *   *   *   *   8
* *   * *   * *   * 9
*   *       *   *   10
\end{lstlisting}
\end{problem}

\begin{problem}
(\tit{Perfect Numbers}) Write a client program \lstinline$PerfectNumbers$ that takes an \lstinline$int$ value $n$ as command-line argument and prints the numbers less than or equal to $n$ that are perfect. Your implementation must use the \lstinline$isPerfect()$ method from \lstinline$Numbers$ to check if a number is perfect.

\begin{lstlisting}[language=bash]
$ java PerfectNumbers 10000
6
28
496
8128
\end{lstlisting}
\end{problem}

\begin{problem}
(\tit{Amicable Pairs}) Write a client program \lstinline$AmicablePairs$ that takes an \lstinline$int$ value $n$ as command-line argument and prints distinct pairs $(i, j)$ with $i \neq j$ and $1 \leq i, j \leq n$, such that $i$ and $j$ are amicable. Your implementation must use the \lstinline$amicable()$ method from \lstinline$Numbers$ to check if two numbers are amicable.

\begin{lstlisting}[language=bash]
$ java AmicablePairs 3000
(220, 284)
(1184, 1210)
(2620, 2924)
\end{lstlisting}
\end{problem}

\noindent \tbf{Files to Submit:}
\begin{enumerate}
\itemsep0em
\item \lstinline$Dragon.java$
\item \lstinline$DragonPlot.java$
\item \lstinline$Numbers.java$
\item \lstinline$PrimeCounter.java$
\item \lstinline$TotativeCounter.java$
\item \lstinline$CoprimePattern.java$
\item \lstinline$PerfectNumbers.java$
\item \lstinline$AmicablePairs.java$
\item \lstinline$report.txt$
\end{enumerate}
\end{document}
